\section{Related Work}
\label{sec:related}

\paragraph{PCM thermal conductivity enhancement.}
The low thermal conductivity of paraffin-based PCMs has motivated extensive research into enhancement strategies. Metal foam infiltration represents one of the most effective approaches: Allen et al.~\cite{allen2015robust} demonstrated that heat pipe-foam-PCM configurations achieved phase change rates 15 times greater than baseline configurations, while Ghalambaz and Zhang~\cite{ghalambaz2020conjugate} investigated conjugate heat transfer in PCM-metal foam heatsinks. Parida et al.~\cite{parida2020effect} specifically examined the role of natural convection within metal foam/PCM composites, finding that convection effects become significant at higher Rayleigh numbers. Liu et al.~\cite{liu2024pore} conducted pore-scale studies revealing the detailed melting mechanisms within open-celled metal foams. Subsequent experimental work has characterized the transient thermal response of metal foam/paraffin composites~\cite{yu2022temperature,duan2021transient}, quantified the effect of orientation on melting and solidification~\cite{allen2015effect}, and compared foam-based strategies with other enhancement methods~\cite{song2015thermal,zhu2016numerical}. Macroscale numerical models for metal foam/paraffin composites have been validated against visualization experiments~\cite{yao2019numerical}, and particle image velocimetry (PIV) measurements have revealed the convection patterns inside partially foam-filled cavities during melting~\cite{liu2022piv}. Graded metal foam composites have recently been characterized experimentally under nonuniform heat flux~\cite{shen2025experimental}.
Extended surfaces offer an alternative enhancement pathway: longitudinal fins in triplex tube heat exchangers~\cite{patel2018thermal,alabidi2013internal}, variable-length fin configurations~\cite{temel2023experimental}, wire-mesh fin designs~\cite{uniyal2024melting}, and finned spherical capsules~\cite{sharma2022effect} have all demonstrated substantial reductions in melting time. The design space of fin geometries for latent heat thermal energy storage (LHTES) systems has been systematically reviewed~\cite{abdullhussein2025review,srivastava2026advanced}, including active techniques that require external power input~\cite{shank2023review}. Recent experimental studies have further explored the coupling of fins with metal foam and helical coils~\cite{tahmasbi2024effects} and the influence of fin configuration on solidification~\cite{fatlawi2024numerical,rao2022thermal}, while fin geometry has been shown to strongly modulate melting rates in double-pipe storage units~\cite{elsayed2025influence}. The melting and solidification behavior of paraffin in shell-and-tube geometries---the workhorse configuration of LHTES research---has been characterized experimentally in both laboratory-scale equipment~\cite{trp2005experimental,gasia2016experimental,ashokkumar2022experimental} and heat-pipe-augmented systems~\cite{motahar2016experimental}, while parametric numerical studies have addressed shell geometry and tube eccentricity~\cite{parsakamkari2024experimental}, dimpled enclosures~\cite{soliman2021melting,soliman2021numerical}, plate-fin storages~\cite{lv2008numerical}, and annular configurations~\cite{kothari2018comprehensive}.
Nanoparticle enhancement, while providing more modest improvements, offers the advantage of minimal structural modification~\cite{murugan2017thermal,bora2026enhanced,bashirpour2021investigation,yadav2020nanoparticle,rudrabhiramu2026nanoparticle,manojkumar2020experimental,farahani2021efficacy}. Nanoparticle-loaded PCMs have also been investigated for cool storage applications~\cite{almeshaal2023numerical}, and carbon-based nano-additives such as quantum dots and MWCNT have been shown to raise the thermal conductivity of paraffin~\cite{emeema2024investigations,sathishkumar2023investigations,george2020novel}. Bio-derived and composite form-stable PCMs further address leakage while improving thermal performance~\cite{yadav2024experimental,gogoi2024thermal}. A distinct class of enhancement employs PCM slurries and microencapsulated suspensions, which exploit latent heat directly in convective flows and have been characterized for both laminar and turbulent duct flows~\cite{charunyakorn1991forced,choi1994forced,alisetti2000forced,inaba2004melting,zhang2003theoretical,delgado2012review}.

\paragraph{PCM-based heat sinks for electronics cooling.}
Passive thermal management of electronics using PCM-filled heat sinks has been studied extensively across a wide range of fin architectures. Plate-fin and multi-fin heat sinks with PCM have been analyzed through three-dimensional transient simulations~\cite{wang2011three}, while hybrid designs combining PCM with both passive and active cooling have been proposed for transient heat load scenarios~\cite{borkar2022transient,marri2023multiple,sevik2025thermal,midhun2025solid}. Experimental studies on circular pin-fin heat sinks demonstrated effective temperature regulation during phase change~\cite{isiam2023experimental,rangarajan2019experimental}, and PCM-based fin heat sinks subjected to transient heat flux shocks showed significant peak temperature reductions~\cite{zhang2025experimental}. Recent work has also combined PCM with graphene nanofluids in hybrid pin-fin architectures~\cite{srevathsan2026experimental}, and explored thermal function improvements of PCM-based systems for electronics cooling~\cite{peng2021thermal}.

\paragraph{PCM in solar thermal systems.}
Solar thermal applications constitute one of the largest application domains of PCM-based storage. Solar air heaters with integrated latent heat storage have been investigated experimentally with various configurations, including finned plates~\cite{kabeel2017improvement,balakrishnan2024experimental}, helical tubes~\cite{saxena2020thermal}, wavy fins~\cite{singh2020numerical}, packed recyclable aluminum cans~\cite{tuncer2023experimental}, and double-pass designs coupled with drying systems~\cite{embiale2023investigation}, and numerical studies have further examined finned absorbers combined with latent heat storage in solar air heaters~\cite{hap2025numerical}. Early foundational studies established the thermal performance of solar air heaters with built-in storage~\cite{fath1995thermal}, and the field has been reviewed comprehensively~\cite{sharmapitchumani2022solar,arkar2015optimization}. Solar cookers with PCM storage have also received attention, from design feasibility studies~\cite{kumar2022design,chauhan2022design} and PCM selection optimization~\cite{anilkumar2021optimum} to recent portable cooker developments~\cite{pinnarat2026development}, and conical cavity receivers with integrated PCM storage have been assessed for parabolic dish collectors~\cite{nandanwar2025performance}. Photovoltaic-thermal (PV/T) collectors integrated with PCM have demonstrated improved electrical and thermal output~\cite{awad2022performance,fakhraei2024experimental}, and PCM-based cold storage for food applications has been experimentally validated~\cite{chinnasamy2021experimental,mane2024experimental,esakkimuthu2013experimental,raj2020transient,famiglietti2021swc,touati2017thermal,shalaby2022experimental}.

\paragraph{PCM in battery and building applications.}
Battery thermal management systems (BTMS) rely increasingly on PCM-based cooling to maintain lithium-ion cells within their safe operating temperature range~\cite{jiang2024battery,lebrouhi2024comparative}. Recent system-level designs combine PCM with liquid cooling~\cite{keyhani2025innovative,pakrouh2023thermal,masthanvali2024battery,wu2026numerical,kumar2025battery}, and nano-enhanced PCMs have been experimentally evaluated for passive battery cooling~\cite{hadawale2025passive}. In the building sector, PCMs are incorporated into envelopes, Trombe walls, and plasters to reduce temperature swings and shift thermal loads~\cite{alyasiri2022energetic,muthuvel2014passive}. Microencapsulated PCM in gypsum plaster and cementitious composites has been shown to improve building thermal storage performance~\cite{wi2021exterior,enteshari2022fabrication,hekimoglu2021thermal}, while Trombe-wall variants containing PCM have been analyzed experimentally and numerically~\cite{zhou2023effect,krason2025thermal,tareemi2026advances}, and room-temperature flexible PCMs with high infrared emissivity have been developed for building energy saving~\cite{kong2024room}.

\paragraph{TPMS geometry and heat transfer.}
The unique geometric properties of TPMS structures---zero mean curvature, high surface area density, and bicontinuous topology---make them inherently attractive for heat transfer applications~\cite{jung2006fluid,tang2023analysis}. Tang et al.~\cite{tang2023analysis} provided a comprehensive performance evaluation of Diamond, Gyroid, and IWP TPMS configurations, establishing quantitative heat transfer correlations. Kaur and Singh~\cite{kaur2021flow} investigated the flow and thermal transport characteristics of Gyroid and Schwarz-P cellular materials, revealing the interplay between lattice geometry and transport properties. Recent work has focused on hybrid and gradient TPMS designs: Chen et al.~\cite{chen2026comprehensive} analyzed Gyroid-Primitive hybrid heat sinks, while Wei et al.~\cite{wei2026convective} proposed novel IWP-Gyroid hybrid structures demonstrating 4--23\% improvements in convective heat transfer coefficient. Barakat and Sun~\cite{barakat2024controlling} explored controlled lattice deformation as a means of further enhancing thermal performance, and gradient TPMS configurations have been evaluated for their thermal performance~\cite{an2026performance}. Convective heat transfer in lattice and TPMS structures for gas-turbine cooling has been assessed numerically, identifying gyroid as the most efficient topology~\cite{casini2024numerical}, while gyroid-type TPMS structures with variable porosity have been analyzed for combined aero engines~\cite{sun2025numerical}. Beyond the thermal domain, TPMS lattices have been evaluated for microprocessor cooling~\cite{ifa2026diamond}, flow boiling in TPMS channels~\cite{hong2026flow,hong2026experimental}, convective performance of sheet-based TPMS~\cite{huo2026convective}, orientation and morphological effects on heat transfer~\cite{son2024effects,son2025effect}, computational case studies of Schwarz-P and gyroid heat sinks~\cite{banas2025evaluation,mian2025computational,nirala2025computational}, acoustic absorption~\cite{godakawela2025sound}, and torsional mechanical response~\cite{skibar2026comparative}. TPMS recuperators for concentrated solar power cycles have also been reviewed, covering hybridized additive manufacturing routes~\cite{letlhare2025hybridized}.

\paragraph{Additive manufacturing of TPMS structures.}
The fabrication of TPMS lattice structures has been greatly enabled by metal additive manufacturing. Catchpole-Smith et al.~\cite{catchpole2019thermal} measured the thermal conductivity of TPMS lattices manufactured via L-PBF, establishing the relationship between relative density and effective conductivity. Chouhan et al.~\cite{chouhan2025heat} experimentally evaluated five L-PBF manufactured TPMS heat sinks (Gyroid, Diamond, Schwarz, Lidinoid, and Split P) against conventional pin-fin designs, demonstrating superior thermal performance. Yanagihara et al.~\cite{yanagihara2026flow} developed flow-priority optimization methods for variable-TPMS lattice heat exchangers. The mechanical properties of additively manufactured aluminum alloy lattices have been extensively characterized: Mazur et al.~\cite{mazur2017mechanical} investigated Ti6Al4V and AlSi12Mg lattice structures, while Noronha et al.~\cite{noronha2025thin} demonstrated highly enhanced specific strength in thin-plate AlSi10Mg lattices, and comparisons of strut-based versus TPMS lattices produced by electron beam melting have quantified the mechanical trade-offs~\cite{sokollu2022mechanical,platek2020investigations}. Stiffness-optimized TPMS infill strategies have further extended the design space of extruded and printed lattices~\cite{wegner2026stiffness}. The thermal conductivity of L-PBF AlSi10Mg itself is process- and orientation-dependent~\cite{elkholy2022characterization,azizi2023process,liu2021inhomogeneous}, motivating alloy design studies for additive manufacturing~\cite{gong2023aluminum,teodorescu2025investigation}. Design and manufacturing guidelines for TPMS structures have been consolidated in recent reference works~\cite{traub2026design}, and reviews of additively manufactured heat exchanger configurations increasingly cover TPMS-based designs~\cite{choudhari2025additive}, while computational design tools based on density-functional theory and generative diffusion models are emerging~\cite{yin2025density,yan2026fourier}.

\paragraph{PCM in structured lattices.}
The integration of PCM with structured lattice or scaffold materials represents the most directly relevant prior work. Piacquadio et al.~\cite{piacquadio2022experimental} provided the first experimental analysis of PCM embedded in additively manufactured lattice structures, demonstrating enhanced thermal energy storage performance. Lebaal et al.~\cite{lebaal2020conjugate} performed conjugate heat transfer analysis of lattice-filled heat exchangers designed for additive manufacturing. Additively manufactured gyroid heat exchangers have meanwhile demonstrated strong thermo-hydraulic performance in both computational and experimental assessments~\cite{mahmoud2023enhancement,yan2024thermohydraulic,lai2025development,yu2025data,palomba2025thermal,daifalla2024multi}. A rapidly growing body of numerical work has examined PCM melting within TPMS lattices, showing accelerated charging in gyroid, diamond, and primitive architectures~\cite{ghafoorian2025numerical,gado2024improving,gado2025thermal,gado2026efficient,su2026thermal,elsihy2026melting,fan2023investigation,avivi2023ihtc,wang2026modulation,molteni2023improving,molteni2026tuning}, and experimental studies have begun to validate these predictions for intermittent heat transfer conditions~\cite{arqam2025computational,wang2024experimental,zhang2023ihtc,zhang2023heat}. However, systematic experimental comparisons of different TPMS topologies (Gyroid vs.\ Primitive vs.\ IWP) under varying heating conditions remain scarce, and the thermal performance of such composites during both melting and solidification has not been comprehensively characterized. The present work addresses this gap by systematically characterizing the thermal response of paraffin-based PCM within AlSiMg TPMS lattices across multiple heating powers and topologies.

\paragraph{Numerical modeling of phase change.}
The enthalpy-porosity method has become the standard numerical approach for simulating solid-liquid phase change~\cite{ebrahimi2019sensitivity,vanriet2024limitations}. This method treats the mushy zone as a pseudo-porous medium, with porosity decreasing from unity (liquid) to zero (solid). The permeability coefficient (mushy zone constant) significantly affects predictions, particularly for non-isothermal phase change~\cite{ebrahimi2019sensitivity}. Natural convection effects during melting have been shown to play a significant role in PCM thermal behavior~\cite{cao2021characterization,souayfane2018melting,bondareva2021natural}, with convection increasing heat loss by more than 10\% at large melt volumes~\cite{bondareva2021natural}. The coupling of wall conduction with convection-dominated melting was established early on~\cite{lacroix1994coupling}, and natural convection at the solid-liquid interface remains an active research topic~\cite{gobin2004natural}. More recent studies have addressed convection-driven melting in inclined and non-uniformly heated cavities~\cite{laouer2021study,bondareva2023natural}, Rayleigh--B{\'e}nard convection during melting under Neumann boundary conditions~\cite{parsazadeh2021numerical}, thermocapillary effects under local heating~\cite{kim2013numerical}, and industrial packed-bed storage systems~\cite{schwarzmayr2023packed,rai2023numerical}. On the experimental side, the thermal conductivity of additively manufactured metal foams has been characterized with an eye toward novel heat exchanger production~\cite{kovalevsky2020experimental}, and additively manufactured lattice heat sinks have been compared against conventional straight-fin designs for railway applications~\cite{batikh2024investigating}.
